% Created 2026-09-21 Mon 12:43
% Intended LaTeX compiler: pdflatex
\documentclass[12pt,a4paper]{article}
\usepackage[utf8]{inputenc}
\usepackage[T1]{fontenc}
\usepackage{graphicx}
\usepackage{grffile}
\usepackage{longtable}
\usepackage{wrapfig}
\usepackage{rotating}
\usepackage[normalem]{ulem}
\usepackage{amsmath}
\usepackage{textcomp}
\usepackage{amssymb}
\usepackage{capt-of}
\usepackage{hyperref}
\usepackage[a4paper,margin=1in]{geometry}
\usepackage{graphicx}
\usepackage{booktabs}
\usepackage{amsmath,amssymb}
\usepackage{float}
\usepackage{hyperref}
\usepackage{xcolor}
\usepackage{listings}
\usepackage{enumitem}
\usepackage{fancyhdr}
\usepackage{longtable}
\hypersetup{colorlinks=true,linkcolor=blue,urlcolor=blue}
\pagestyle{fancy}
\fancyhf{}
\lhead{ICS1512}
\rhead{Machine Learning Algorithms Laboratory}
\cfoot{\thepage}
\setlength{\headheight}{14.5pt}
\lstset{language=Python,basicstyle=\ttfamily\small,numbers=left,frame=single,breaklines=true,showstringspaces=false}
\author{Emmanuel Maximus}
\date{\today}
\title{}
\hypersetup{
 pdfauthor={Emmanuel Maximus},
 pdftitle={},
 pdfkeywords={},
 pdfsubject={},
 pdfcreator={Emacs 27.1 (Org mode 9.3)}, 
 pdflang={English}}
\begin{document}

\tableofcontents

l\#+OPTIONS: title:nil toc:nil
\begin{center}
\begin{tabular}{ll}
\hline
Experiment No. & 8\\
Experiment Title & Clustering Human Activity Recognition Data using K-Means, DBSCAN, and Hierarchical Clustering\\
Student Name & Emmanuel Maximus Ravichandran\\
Register Number & 3122247001015\\
Faculty & Dr. Poreddy Ajay Kumar Reddy\\
Submission Date & \\
\hline
\end{tabular}
\end{center}

\section{Objective}
\label{sec:org8ca33be}

To implement and analyze the performance of three clustering algorithms —
K-Means, DBSCAN, and Hierarchical Agglomerative Clustering (Ward linkage) —
on the Human Activity Recognition (HAR) Using Smartphones dataset, visualize
the clusters formed by each, and compare them against the six ground-truth
activity labels.

\section{Problem Statement}
\label{sec:orgc0a9f2a}

\textbf{Input:} 561 time- and frequency-domain features extracted from smartphone
accelerometer and gyroscope signals (50 Hz, windowed into 2.56s segments)
for 30 volunteers performing 6 activities.

\textbf{Output:} Cluster assignments produced independently by K-Means, DBSCAN, and
Hierarchical (Ward) clustering, with no access to the true activity label
during clustering.

\textbf{Objective:} Since clustering is unsupervised, the true activity labels
(WALKING, WALKING$\backslash$\textsubscript{UPSTAIRS}, WALKING$\backslash$\textsubscript{DOWNSTAIRS}, SITTING, STANDING, LAYING)
are used only afterwards, to externally validate how well each algorithm's
clusters correspond to real human activities.

\section{Dataset Description}
\label{sec:org1150110}

\subsection{Setup and Imports}
\label{sec:org4fea716}

\begin{verbatim}
import warnings
warnings.filterwarnings("ignore")

import time, json
import numpy as np
import pandas as pd
import matplotlib
matplotlib.use("Agg")
import matplotlib.pyplot as plt
import seaborn as sns

from sklearn.preprocessing import StandardScaler
from sklearn.decomposition import PCA
from sklearn.manifold import TSNE
from sklearn.cluster import KMeans, DBSCAN
from sklearn.neighbors import NearestNeighbors
from sklearn.metrics import (silhouette_score, davies_bouldin_score, calinski_harabasz_score,
			      adjusted_rand_score, normalized_mutual_info_score, confusion_matrix)
from scipy.cluster.hierarchy import linkage, dendrogram, fcluster
from scipy.stats import mode

np.random.seed(42)
sns.set_style("whitegrid")
print("Environment ready.")
\end{verbatim}

\begin{verbatim}
Python 3.10.12 (main, Aug 31 2026, 10:18:17) [GCC 11.4.0] on linux
Type "help", "copyright", "credits" or "license" for more information.
Environment ready.
python.el: native completion setup loaded
\end{verbatim}

\subsection{Loading the Dataset}
\label{sec:org77af6e3}

\begin{verbatim}
BASE = "/home/claude/har/UCI HAR Dataset"

features = pd.read_csv(f"{BASE}/features.txt", sep=r"\s+", header=None, names=["idx", "feature"])
feat_names = features["feature"].tolist()

seen, uniq = {}, []
for f in feat_names:
    if f in seen:
	seen[f] += 1
	uniq.append(f"{f}_{seen[f]}")
    else:
	seen[f] = 0
	uniq.append(f)

X_train = pd.read_csv(f"{BASE}/train/X_train.txt", sep=r"\s+", header=None, names=uniq)
X_test  = pd.read_csv(f"{BASE}/test/X_test.txt",  sep=r"\s+", header=None, names=uniq)
y_train = pd.read_csv(f"{BASE}/train/y_train.txt", header=None, names=["Activity"])
y_test  = pd.read_csv(f"{BASE}/test/y_test.txt",  header=None, names=["Activity"])

act_labels = pd.read_csv(f"{BASE}/activity_labels.txt", sep=r"\s+", header=None, names=["id", "name"])
act_map = dict(zip(act_labels["id"], act_labels["name"]))
act_order = ["WALKING", "WALKING_UPSTAIRS", "WALKING_DOWNSTAIRS", "SITTING", "STANDING", "LAYING"]

X = pd.concat([X_train, X_test], axis=0).reset_index(drop=True)
y = pd.concat([y_train, y_test], axis=0).reset_index(drop=True)
y_named = y["Activity"].map(act_map)
y_code = y_named.map({a: i for i, a in enumerate(act_order)}).values

print("Feature matrix shape:", X.shape)
print("Duplicate raw feature names fixed:", sum(v > 0 for v in seen.values()))
print("Missing values in X:", int(X.isna().sum().sum()))
print(y_named.value_counts())
\end{verbatim}

\begin{verbatim}
Traceback (most recent call last):
  File "<stdin>", line 1, in <module>
  File "/tmp/babel-cBWeqx/python-HNlEOF", line 3, in <module>
    features = pd.read_csv(f"{BASE}/features.txt", sep=r"\s+", header=None, names=["idx", "feature"])
  File "/usr/lib/python3/dist-packages/pandas/util/_decorators.py", line 311, in wrapper
    return func(*args, **kwargs)
  File "/usr/lib/python3/dist-packages/pandas/io/parsers/readers.py", line 586, in read_csv
    return _read(filepath_or_buffer, kwds)
  File "/usr/lib/python3/dist-packages/pandas/io/parsers/readers.py", line 482, in _read
    parser = TextFileReader(filepath_or_buffer, **kwds)
  File "/usr/lib/python3/dist-packages/pandas/io/parsers/readers.py", line 811, in __init__
    self._engine = self._make_engine(self.engine)
  File "/usr/lib/python3/dist-packages/pandas/io/parsers/readers.py", line 1040, in _make_engine
    return mapping[engine](self.f, **self.options)  # type: ignore[call-arg]
  File "/usr/lib/python3/dist-packages/pandas/io/parsers/c_parser_wrapper.py", line 51, in __init__
    self._open_handles(src, kwds)
  File "/usr/lib/python3/dist-packages/pandas/io/parsers/base_parser.py", line 222, in _open_handles
    self.handles = get_handle(
  File "/usr/lib/python3/dist-packages/pandas/io/common.py", line 702, in get_handle
    handle = open(
FileNotFoundError: [Errno 2] No such file or directory: '/home/claude/har/UCI HAR Dataset/features.txt'
\end{verbatim}

Dataset Description table (values confirmed by the code above):

\begin{center}
\begin{tabular}{lr}
\hline
Dataset Name & Human Activity Recognition Using Smartphones\\
Dataset Source & UCI Machine Learning Repository\\
Number of Samples & 10299\\
Number of Features & 561\\
Number of Classes & 6\\
Missing Values & 0\\
Train-Test Split & Original UCI split (7352/2947), recombined for clustering\\
\hline
\end{tabular}
\end{center}

\section{Exploratory Data Analysis}
\label{sec:org85d38f2}

Include all relevant visualizations.
\begin{itemize}
\item Dataset overview
\item Statistical summary
\item Missing value analysis
\item Class distribution
\item Correlation matrix
\item Feature distribution
\end{itemize}

\begin{verbatim}
print(X.describe().T[["mean", "std", "min", "max"]].head(8))
\end{verbatim}

\begin{verbatim}
Traceback (most recent call last):
  File "<stdin>", line 1, in <module>
NameError: name 'X' is not defined
\end{verbatim}


\begin{verbatim}
na_counts = X.isna().sum()
print("Total missing cells:", int(na_counts.sum()))
print("Columns with any missing value:", int((na_counts > 0).sum()))
\end{verbatim}

\begin{verbatim}
Traceback (most recent call last):
  File "<stdin>", line 1, in <module>
  File "/tmp/babel-cBWeqx/python-4t4Nhs", line 1, in <module>
    na_counts = X.isna().sum()
NameError: name 'X' is not defined
\end{verbatim}


\begin{verbatim}
plt.figure(figsize=(7, 4.5))
sns.countplot(y=y_named, order=act_order, palette="viridis")
plt.title("Class Distribution of Activities")
plt.xlabel("Count"); plt.ylabel("Activity")
plt.tight_layout()
plt.savefig("/home/claude/har/nbassets/class_distribution.png", dpi=120)
plt.close()
"/home/claude/har/nbassets/class_distribution.png"
\end{verbatim}
\captionof{figure}{\label{org0978e3b}Class distribution of the six activities}

\url{Traceback (most recent call last): File "<stdin>", line 1, in <module> File "/tmp/babel-cBWeqx/python-FzQIah", line 2, in <module> sns.countplot(y=y\_named, order=act\_order, palette="viridis") NameError: name 'y\_named' is not defined. Did you mean: '\_\_name\_\_'?}

\textbf{Inference:} The class distribution is close to balanced across the six
activities (1406-1944 samples each), so no resampling is needed before
clustering.

\begin{verbatim}
corr_cols = X.columns[:40]
plt.figure(figsize=(9, 7))
sns.heatmap(X[corr_cols].corr(), cmap="coolwarm", center=0, xticklabels=False, yticklabels=False)
plt.title("Correlation Matrix (first 40 features)")
plt.tight_layout()
plt.savefig("/home/claude/har/nbassets/correlation_matrix.png", dpi=120)
plt.close()
"/home/claude/har/nbassets/correlation_matrix.png"
\end{verbatim}
\captionof{figure}{\label{org08bc1f9}Correlation matrix of the first 40 features}

\url{Traceback (most recent call last): File "<stdin>", line 1, in <module> File "/tmp/babel-cBWeqx/python-W6Lab8", line 1, in <module> corr\_cols = X.columns[:40] NameError: name 'X' is not defined}

\textbf{Inference:} Large contiguous blocks of highly correlated features appear,
since many features are simple mean/std/energy statistics of the same
handful of underlying accelerometer/gyroscope signal axes. This heavy
redundancy is exactly what makes the PCA dimensionality reduction step
below so effective.

\begin{verbatim}
feat_pick = ["tBodyAcc-mean()-X", "tGravityAcc-mean()-X", "tBodyAccMag-mean()", "fBodyAccMag-mean()"]
fig, axes = plt.subplots(2, 2, figsize=(10, 7))
for ax, col in zip(axes.ravel(), feat_pick):
    for act in act_order:
	sns.kdeplot(X.loc[y_named == act, col], ax=ax, label=act)
    ax.set_title(col)
axes[0, 0].legend(fontsize=7)
plt.tight_layout()
plt.savefig("/home/claude/har/nbassets/feature_distributions.png", dpi=120)
plt.close()
"/home/claude/har/nbassets/feature_distributions.png"
\end{verbatim}
\captionof{figure}{\label{org54d49f2}Distributions of four representative features, split by activity}

\url{Traceback (most recent call last): File "<stdin>", line 1, in <module> File "/tmp/babel-cBWeqx/python-PhIz40", line 4, in <module> for act in act\_order: NameError: name 'act\_order' is not defined}

\textbf{Inference:} Static activities (SITTING, STANDING, LAYING) separate clearly
from dynamic activities (the WALKING variants) along magnitude-based
features, but SITTING and STANDING largely overlap with each other, and the
three WALKING sub-types overlap heavily with one another. This foreshadows
the clustering difficulty seen later, where these pairs are the hardest to
tell apart.

\section{Data Preprocessing}
\label{sec:orga344264}

\begin{itemize}
\item Missing value handling: not required -- 0 missing values in the dataset.
\item Encoding: activity labels are integer-coded 1-6 in the raw files; mapped
to their names via \texttt{activity\_labels.txt}, and kept aside only for
post-hoc external validation.
\item Feature scaling: \texttt{StandardScaler} (zero mean, unit variance) applied to
all 561 features.
\item Train-test split: not applicable in the supervised sense -- the original
split is recombined into one set for clustering.
\item Feature engineering: PCA is applied after scaling, to reduce
dimensionality for clustering and to produce 2D projections for
visualization.
\end{itemize}

\begin{verbatim}
scaler = StandardScaler()
Xs = scaler.fit_transform(X)

pca_full = PCA(random_state=42).fit(Xs)
cum_var = np.cumsum(pca_full.explained_variance_ratio_)
n90 = int(np.argmax(cum_var >= 0.90) + 1)
n95 = int(np.argmax(cum_var >= 0.95) + 1)

plt.figure(figsize=(7, 4.5))
plt.plot(range(1, len(cum_var) + 1), cum_var, lw=2)
plt.axhline(0.90, color="r", ls="--", label="90% variance")
plt.axvline(n90, color="r", ls=":")
plt.xlabel("Number of Principal Components"); plt.ylabel("Cumulative Explained Variance")
plt.title("PCA Cumulative Explained Variance")
plt.legend()
plt.tight_layout()
plt.savefig("/home/claude/har/nbassets/pca_variance.png", dpi=120)
plt.close()
"/home/claude/har/nbassets/pca_variance.png"
\end{verbatim}
\captionof{figure}{\label{org2c4a213}PCA cumulative explained variance}

\url{Traceback (most recent call last): File "<stdin>", line 1, in <module> File "/tmp/babel-cBWeqx/python-sYOLdW", line 2, in <module> Xs = scaler.fit\_transform(X) NameError: name 'X' is not defined}

\begin{verbatim}
print(f"Components needed for 90% variance: {n90}")
print(f"Components needed for 95% variance: {n95}")

pca90 = PCA(n_components=n90, random_state=42)
Xp = pca90.fit_transform(Xs)
pca2 = PCA(n_components=2, random_state=42)
X2 = pca2.fit_transform(Xs)
print("Xp (PCA-90%) shape:", Xp.shape, " X2 (PCA-2D) shape:", X2.shape)
print("PC1+PC2 explained variance:", round(sum(pca2.explained_variance_ratio_), 4))
\end{verbatim}

\begin{verbatim}
Traceback (most recent call last):
  File "<stdin>", line 1, in <module>
  File "/tmp/babel-cBWeqx/python-ZhXzfX", line 1, in <module>
    print(f"Components needed for 90% variance: {n90}")
NameError: name 'n90' is not defined
\end{verbatim}


\begin{verbatim}
t0 = time.time()
Xtsne = TSNE(n_components=2, random_state=42, perplexity=30, init="pca").fit_transform(Xp)
print(f"t-SNE fit time: {time.time()-t0:.1f}s, output shape: {Xtsne.shape}")
\end{verbatim}

\begin{verbatim}
Traceback (most recent call last):
  File "<stdin>", line 1, in <module>
  File "/tmp/babel-cBWeqx/python-uonYR0", line 2, in <module>
    Xtsne = TSNE(n_components=2, random_state=42, perplexity=30, init="pca").fit_transform(Xp)
NameError: name 'Xp' is not defined. Did you mean: 'np'?
\end{verbatim}


\begin{verbatim}
fig, axes = plt.subplots(1, 2, figsize=(13, 5.5))
palette = sns.color_palette("tab10", 6)
for i, act in enumerate(act_order):
    m = (y_named == act).values
    axes[0].scatter(X2[m, 0], X2[m, 1], s=5, color=palette[i], label=act, alpha=0.6)
    axes[1].scatter(Xtsne[m, 0], Xtsne[m, 1], s=5, color=palette[i], label=act, alpha=0.6)
axes[0].set_title("PCA (2D) - Ground Truth Activities"); axes[0].set_xlabel("PC1"); axes[0].set_ylabel("PC2")
axes[1].set_title("t-SNE (2D) - Ground Truth Activities"); axes[1].set_xlabel("t-SNE1"); axes[1].set_ylabel("t-SNE2")
axes[0].legend(fontsize=7, markerscale=2)
plt.tight_layout()
plt.savefig("/home/claude/har/nbassets/ground_truth_pca_tsne.png", dpi=120)
plt.close()
"/home/claude/har/nbassets/ground_truth_pca_tsne.png"
\end{verbatim}
\captionof{figure}{\label{org3138381}Ground-truth activities in PCA(2D) and t-SNE(2D)}

\url{Traceback (most recent call last): File "<stdin>", line 1, in <module> File "/tmp/babel-cBWeqx/python-tKP1E5", line 3, in <module> for i, act in enumerate(act\_order): NameError: name 'act\_order' is not defined}

\textbf{Inference:} t-SNE gives a much cleaner separation of the six activities
than linear PCA -- PCA's first two components only capture about 57\% of
total variance and mostly separate dynamic from static activities, while
SITTING/STANDING/LAYING and the three WALKING variants stay visibly
clustered together within their own broad regions even in 2D PCA.

\section{Machine Learning Algorithm}
\label{sec:orge38bec0}

\subsection{K-Means Clustering}
\label{sec:org1796a4b}
\begin{itemize}
\item Working principle: partitions data into \(k\) clusters by iteratively
assigning each point to its nearest centroid (Euclidean distance) and
recomputing centroids as the mean of assigned points, until convergence.
\item Mathematical formulation: minimizes
\(\text{WCSS} = \sum_{j=1}^{k}\sum_{x_i \in C_j} \lVert x_i - \mu_j \rVert^2\).
\item Advantages: simple, fast, scales well; guaranteed to converge.
\item Limitations: requires \(k\) in advance; assumes roughly spherical,
similarly-sized clusters; sensitive to initialization and feature scale.
\end{itemize}

\subsection{DBSCAN}
\label{sec:org6c217f1}
\begin{itemize}
\item Working principle: groups points that are densely packed (density-
connected within radius \texttt{eps}, with at least \texttt{minPts} neighbors), and
marks points in low-density regions as noise.
\item Mathematical formulation: a point \(p\) is a \textbf{core point} if
\(|N_{eps}(p)| \geq \text{minPts}\); clusters are maximal sets of
density-connected points.
\item Advantages: no need to fix \(k\) in advance; finds arbitrarily-shaped
clusters; naturally identifies outliers.
\item Limitations: sensitive to \texttt{eps=/=minPts}; struggles in high-dimensional
spaces where distances concentrate (demonstrated in the Hyperparameter
Selection section below).
\end{itemize}

\subsection{Hierarchical Agglomerative Clustering (Ward linkage)}
\label{sec:org34f5004}
\begin{itemize}
\item Working principle: starts with every point as its own cluster and
repeatedly merges the pair of clusters whose merger causes the smallest
increase in total within-cluster variance, building a dendrogram.
\item Mathematical formulation: Ward's criterion merges clusters \(A, B\)
minimizing \(\Delta(A,B) = \frac{|A||B|}{|A|+|B|}\lVert \mu_A - \mu_B \rVert^2\).
\item Advantages: no need to pre-specify \(k\); produces a full hierarchy;
deterministic.
\item Limitations: \(O(n^2)\) or worse time/memory; greedy merges cannot be
undone; favors compact, similarly-sized clusters.
\end{itemize}

\section{Experimental Setup}
\label{sec:org358776b}

\begin{verbatim}
import sklearn, scipy, platform
print(f"Python Version: {platform.python_version()}")
print(f"Scikit-Learn Version: {sklearn.__version__}")
print(f"SciPy Version: {scipy.__version__}")
print(f"Pandas Version: {pd.__version__}")
print(f"Operating System: {platform.system()} {platform.release()}")
print(f"Machine: {platform.machine()}")
\end{verbatim}

\begin{verbatim}
Python Version: 3.10.12
Scikit-Learn Version: 1.7.2
SciPy Version: 1.8.0
Pandas Version: 1.3.5
Operating System: Linux 6.8.0-138-generic
Machine: x86_64
\end{verbatim}


(IDE and hardware depend on where this file is ultimately executed --
record accordingly; the environment above is the sandbox where this
analysis was authored and fully executed.)

\section{Hyperparameter Selection}
\label{sec:orgcc600af}

\subsection{K-Means -- Elbow Method}
\label{sec:org5321432}

\begin{verbatim}
ks = list(range(2, 9))
wcss, sil = [], []
for k in ks:
    km = KMeans(n_clusters=k, n_init=10, random_state=42).fit(Xp)
    wcss.append(km.inertia_)
    sil.append(silhouette_score(Xp, km.labels_, sample_size=4000, random_state=42))

elbow_df = pd.DataFrame({"k": ks, "WCSS (Inertia)": wcss, "Silhouette Score": sil})

fig, axes = plt.subplots(1, 2, figsize=(12, 4.5))
axes[0].plot(ks, wcss, marker="o"); axes[0].set_xlabel("k"); axes[0].set_ylabel("WCSS (Inertia)")
axes[0].set_title("Elbow Curve: k vs WCSS")
axes[1].plot(ks, sil, marker="o", color="darkorange"); axes[1].set_xlabel("k"); axes[1].set_ylabel("Silhouette Score")
axes[1].set_title("k vs Silhouette Score")
plt.tight_layout()
plt.savefig("/home/claude/har/nbassets/kmeans_elbow_silhouette.png", dpi=120)
plt.close()
"/home/claude/har/nbassets/kmeans_elbow_silhouette.png"
\end{verbatim}
\captionof{figure}{\label{orgb4dbf76}Elbow curve (k vs WCSS) and k vs Silhouette Score}

\url{Traceback (most recent call last): File "<stdin>", line 1, in <module> File "/tmp/babel-cBWeqx/python-a6cbhe", line 4, in <module> km = KMeans(n\_clusters=k, n\_init=10, random\_state=42).fit(Xp) NameError: name 'Xp' is not defined. Did you mean: 'np'?}

\begin{verbatim}
print(elbow_df.to_string(index=False))
best_k_internal = ks[int(np.argmax(sil))]
print(f"\nBest k by internal Silhouette score: {best_k_internal}")
\end{verbatim}

\begin{verbatim}
Traceback (most recent call last):
  File "<stdin>", line 1, in <module>
  File "/tmp/babel-cBWeqx/python-GrCRqh", line 1, in <module>
    print(elbow_df.to_string(index=False))
NameError: name 'elbow_df' is not defined
\end{verbatim}


\textbf{Selected k:} internal metrics point to \(k=2\) (moving vs. not moving).
Since the goal is also to compare against the 6 known activities, \(k=6\) is
used as the primary model below, with \(k=2\) kept as a secondary comparison
(see Discussion).

\begin{verbatim}
km_final = KMeans(n_clusters=6, n_init=10, random_state=42).fit(Xp)
km_k2 = KMeans(n_clusters=best_k_internal, n_init=10, random_state=42).fit(Xp)
labels_km, labels_km2 = km_final.labels_, km_k2.labels_
print("K-Means (k=6) cluster sizes:", np.bincount(labels_km))
print("K-Means (k=2) cluster sizes:", np.bincount(labels_km2))
\end{verbatim}

\begin{verbatim}
Traceback (most recent call last):
  File "<stdin>", line 1, in <module>
  File "/tmp/babel-cBWeqx/python-qYt7Tm", line 1, in <module>
    km_final = KMeans(n_clusters=6, n_init=10, random_state=42).fit(Xp)
NameError: name 'Xp' is not defined. Did you mean: 'np'?
\end{verbatim}

\subsection{DBSCAN -- eps and minPts Tuning}
\label{sec:org2c95102}

\begin{verbatim}
minPts = 10
nn = NearestNeighbors(n_neighbors=minPts).fit(X2)
dist, _kdist_idx = nn.kneighbors(X2)
kdist = np.sort(dist[:, -1])

plt.figure(figsize=(7, 4.5))
plt.plot(kdist)
plt.axhline(1.0, color="r", ls="--", label="chosen eps = 1.0")
plt.xlabel("Points sorted by distance"); plt.ylabel(f"Distance to {minPts}-th nearest neighbor")
plt.title("k-distance Graph (for eps selection)")
plt.legend()
plt.tight_layout()
plt.savefig("/home/claude/har/nbassets/dbscan_kdistance.png", dpi=120)
plt.close()
"/home/claude/har/nbassets/dbscan_kdistance.png"
\end{verbatim}
\captionof{figure}{\label{org5ad557c}k-distance graph used to select DBSCAN's eps}

\url{Traceback (most recent call last): File "<stdin>", line 1, in <module> File "/tmp/babel-cBWeqx/python-aBIxju", line 2, in <module> nn = NearestNeighbors(n\_neighbors=minPts).fit(X2) NameError: name 'X2' is not defined}

\begin{verbatim}
eps_sweep = [0.3, 0.4, 0.5, 0.6, 0.7, 0.8, 1.0, 1.2, 1.5]
rows = []
for eps in eps_sweep:
    db = DBSCAN(eps=eps, min_samples=minPts).fit(X2)
    labs = db.labels_
    n_clusters = len(set(labs)) - (1 if -1 in labs else 0)
    noise = int((labs == -1).sum())
    rows.append((eps, n_clusters, noise, round(noise / len(labs) * 100, 1)))

sweep_df = pd.DataFrame(rows, columns=["eps", "Clusters", "Noise Points", "Noise %"])
print(sweep_df.to_string(index=False))
\end{verbatim}

\begin{verbatim}
Traceback (most recent call last):
  File "<stdin>", line 1, in <module>
  File "/tmp/babel-cBWeqx/python-GOt5fE", line 4, in <module>
    db = DBSCAN(eps=eps, min_samples=minPts).fit(X2)
NameError: name 'X2' is not defined
\end{verbatim}


\textbf{Chosen parameters: eps = 1.0, minPts = 10} -- the smallest eps giving
exactly 6 clusters (matching the number of true activities) with noise
held under 10\%.

Curse-of-dimensionality check -- the same search attempted directly on the
65-dimensional PCA space (minPts = 2x65 = 130, the standard heuristic)
essentially fails:

\begin{verbatim}
db_highdim = DBSCAN(eps=16, min_samples=2 * n90).fit(Xp)
n_hd = len(set(db_highdim.labels_)) - (1 if -1 in db_highdim.labels_ else 0)
noise_hd = int((db_highdim.labels_ == -1).sum())
print(f"DBSCAN on 65-D PCA space (eps=16, minPts={2*n90}): {n_hd} cluster(s), {noise_hd}/{len(db_highdim.labels_)} noise ({noise_hd/len(db_highdim.labels_)*100:.1f}%)")
print("-> distances concentrate in high dimensions, so density-based clustering is applied on the 2D PCA projection instead.")
\end{verbatim}

\begin{verbatim}
Traceback (most recent call last):
  File "<stdin>", line 1, in <module>
  File "/tmp/babel-cBWeqx/python-MuMt4P", line 1, in <module>
    db_highdim = DBSCAN(eps=16, min_samples=2 * n90).fit(Xp)
NameError: name 'n90' is not defined
\end{verbatim}


\begin{verbatim}
db_final = DBSCAN(eps=1.0, min_samples=minPts).fit(X2)
labels_db = db_final.labels_
n_db = len(set(labels_db)) - (1 if -1 in labels_db else 0)
print(f"Final DBSCAN: {n_db} clusters, {int((labels_db==-1).sum())} noise points ({(labels_db==-1).mean()*100:.1f}%)")
print("Cluster sizes:", {int(c): int((labels_db == c).sum()) for c in sorted(set(labels_db))})
\end{verbatim}

\begin{verbatim}
Traceback (most recent call last):
  File "<stdin>", line 1, in <module>
  File "/tmp/babel-cBWeqx/python-iGEnM7", line 1, in <module>
    db_final = DBSCAN(eps=1.0, min_samples=minPts).fit(X2)
NameError: name 'X2' is not defined
\end{verbatim}

\subsection{Hierarchical Agglomerative Clustering -- Ward Linkage}
\label{sec:orgab2926f}

\begin{verbatim}
Z = linkage(Xp, method="ward")

plt.figure(figsize=(10, 5.5))
dendrogram(Z, truncate_mode="lastp", p=30, leaf_rotation=90.0, leaf_font_size=8, show_contracted=True)
plt.title("Dendrogram (Ward Linkage, truncated to last 30 merges)")
plt.xlabel("Cluster size (or sample index)"); plt.ylabel("Ward Distance")
plt.tight_layout()
plt.savefig("/home/claude/har/nbassets/dendrogram.png", dpi=120)
plt.close()
"/home/claude/har/nbassets/dendrogram.png"
\end{verbatim}
\captionof{figure}{\label{orgfbe45cb}Dendrogram (Ward linkage, truncated to last 30 merges)}

\url{Traceback (most recent call last): File "<stdin>", line 1, in <module> File "/tmp/babel-cBWeqx/python-IY7DVr", line 1, in <module> Z = linkage(Xp, method="ward") NameError: name 'Xp' is not defined. Did you mean: 'np'?}

\begin{verbatim}
labels_hac = fcluster(Z, t=6, criterion="maxclust") - 1
print("Hierarchical (Ward, k=6) cluster sizes:", np.bincount(labels_hac))
\end{verbatim}

\begin{verbatim}
Traceback (most recent call last):
  File "<stdin>", line 1, in <module>
  File "/tmp/babel-cBWeqx/python-aCHQhN", line 1, in <module>
    labels_hac = fcluster(Z, t=6, criterion="maxclust") - 1
NameError: name 'Z' is not defined
\end{verbatim}

\section{Results}
\label{sec:org074b74a}

\begin{verbatim}
def internal_metrics(Xspace, labels, exclude_noise=False):
    if exclude_noise:
	mask = labels != -1
	Xu, lu = Xspace[mask], labels[mask]
    else:
	Xu, lu = Xspace, labels
    return dict(
	Silhouette=silhouette_score(Xu, lu, sample_size=4000, random_state=42),
	Davies_Bouldin=davies_bouldin_score(Xu, lu),
	Calinski_Harabasz=calinski_harabasz_score(Xu, lu)
    )

def external_metrics(true, pred):
    return dict(ARI=adjusted_rand_score(true, pred), NMI=normalized_mutual_info_score(true, pred))

metrics = {}
metrics["K-Means (k=6)"] = {**internal_metrics(Xp, labels_km), **external_metrics(y_code, labels_km)}
metrics["DBSCAN"] = {**internal_metrics(X2, labels_db, exclude_noise=True), **external_metrics(y_code, labels_db)}
metrics["Hierarchical (Ward)"] = {**internal_metrics(Xp, labels_hac), **external_metrics(y_code, labels_hac)}

metrics_df = pd.DataFrame(metrics).T
print(metrics_df.round(4).to_string())
\end{verbatim}

\begin{verbatim}
Traceback (most recent call last):
  File "<stdin>", line 1, in <module>
  File "/tmp/babel-cBWeqx/python-UIpaKb", line 17, in <module>
    metrics["K-Means (k=6)"] = {**internal_metrics(Xp, labels_km), **external_metrics(y_code, labels_km)}
NameError: name 'Xp' is not defined. Did you mean: 'np'?
\end{verbatim}

\section{Result Visualization}
\label{sec:org97e001f}

Include all graphs generated during the experiment; each with figure,
caption, and inference below.

\begin{verbatim}
fig, axes = plt.subplots(1, 2, figsize=(13, 5.5))
axes[0].scatter(X2[:, 0], X2[:, 1], c=labels_km, cmap="tab10", s=5, alpha=0.6)
axes[0].set_title("K-Means (k=6) clusters - PCA(2D)")
axes[1].scatter(Xtsne[:, 0], Xtsne[:, 1], c=labels_km, cmap="tab10", s=5, alpha=0.6)
axes[1].set_title("K-Means (k=6) clusters - t-SNE(2D)")
plt.tight_layout()
plt.savefig("/home/claude/har/nbassets/kmeans_clusters.png", dpi=120)
plt.close()
"/home/claude/har/nbassets/kmeans_clusters.png"
\end{verbatim}
\captionof{figure}{\label{org48e4d9b}K-Means (k=6) clusters in PCA(2D) and t-SNE(2D)}

\url{Traceback (most recent call last): File "<stdin>", line 1, in <module> File "/tmp/babel-cBWeqx/python-ojpnXF", line 2, in <module> axes[0].scatter(X2[:, 0], X2[:, 1], c=labels\_km, cmap="tab10", s=5, alpha=0.6) NameError: name 'X2' is not defined}

\textbf{Inference:} K-Means (k=6) cleanly isolates LAYING but splits the dynamic-
activity region into clusters that don't map 1:1 onto the WALKING
variants, and merges SITTING with STANDING.

\begin{verbatim}
uniq_labels = sorted(set(labels_db))
cmap = plt.cm.tab10
plt.figure(figsize=(6.5, 5.5))
for u in uniq_labels:
    m = labels_db == u
    if u == -1:
	plt.scatter(X2[m, 0], X2[m, 1], c="lightgray", s=5, label="Noise", alpha=0.5)
    else:
	plt.scatter(X2[m, 0], X2[m, 1], color=cmap(u % 10), s=5, label=f"Cluster {u}", alpha=0.7)
plt.legend(fontsize=7, markerscale=2)
plt.title("DBSCAN (eps=1.0, minPts=10) clusters - PCA(2D)")
plt.tight_layout()
plt.savefig("/home/claude/har/nbassets/dbscan_clusters.png", dpi=120)
plt.close()
"/home/claude/har/nbassets/dbscan_clusters.png"
\end{verbatim}
\captionof{figure}{\label{orgfb71483}DBSCAN (eps=1.0, minPts=10) clusters in PCA(2D)}

\url{Traceback (most recent call last): File "<stdin>", line 1, in <module> File "/tmp/babel-cBWeqx/python-ocSuIc", line 1, in <module> uniq\_labels = sorted(set(labels\_db)) NameError: name 'labels\_db' is not defined}

\textbf{Inference:} DBSCAN finds two large, well-separated density regions
(roughly dynamic vs. static) plus a few small tight sub-clusters and a
thin scatter of noise along the boundary.

\begin{verbatim}
fig, axes = plt.subplots(1, 2, figsize=(13, 5.5))
axes[0].scatter(X2[:, 0], X2[:, 1], c=labels_hac, cmap="tab10", s=5, alpha=0.6)
axes[0].set_title("Hierarchical (Ward, k=6) clusters - PCA(2D)")
axes[1].scatter(Xtsne[:, 0], Xtsne[:, 1], c=labels_hac, cmap="tab10", s=5, alpha=0.6)
axes[1].set_title("Hierarchical (Ward, k=6) clusters - t-SNE(2D)")
plt.tight_layout()
plt.savefig("/home/claude/har/nbassets/hac_clusters.png", dpi=120)
plt.close()
"/home/claude/har/nbassets/hac_clusters.png"
\end{verbatim}
\captionof{figure}{\label{org7a835d9}Hierarchical (Ward, k=6) clusters in PCA(2D) and t-SNE(2D)}

\url{Traceback (most recent call last): File "<stdin>", line 1, in <module> File "/tmp/babel-cBWeqx/python-8INwrK", line 2, in <module> axes[0].scatter(X2[:, 0], X2[:, 1], c=labels\_hac, cmap="tab10", s=5, alpha=0.6) NameError: name 'X2' is not defined}

\textbf{Inference:} Ward linkage produces a visibly cleaner LAYING vs. STANDING
split than K-Means, matching its higher ARI/NMI in the Results table.

\begin{verbatim}
def cluster_to_activity_confusion(pred, exclude_noise_label=None):
    pred = np.array(pred)
    mapping = {}
    for c in set(pred):
	if exclude_noise_label is not None and c == exclude_noise_label:
	    continue
	mask = pred == c
	maj = mode(y_code[mask], keepdims=False).mode
	mapping[c] = maj
    mapped = np.array([mapping.get(c, -1) for c in pred])
    valid = mapped != -1
    cm = confusion_matrix(y_code[valid], mapped[valid], labels=list(range(6)))
    return cm, mapping

cm_km, map_km = cluster_to_activity_confusion(labels_km)
cm_db, map_db = cluster_to_activity_confusion(labels_db, exclude_noise_label=-1)
cm_hac, map_hac = cluster_to_activity_confusion(labels_hac)

fig, axes = plt.subplots(1, 3, figsize=(16, 5))
for ax, cm, name in zip(axes, [cm_km, cm_db, cm_hac], ["K-Means", "DBSCAN", "Hierarchical"]):
    sns.heatmap(cm, annot=True, fmt="d", cmap="Blues", xticklabels=act_order, yticklabels=act_order, ax=ax, cbar=False)
    ax.set_title(f"{name} (mapped to activity)")
    ax.set_xlabel("Predicted (majority-vote)"); ax.set_ylabel("True")
    plt.setp(ax.get_xticklabels(), rotation=45, ha="right")
plt.tight_layout()
plt.savefig("/home/claude/har/nbassets/confusion_matrices.png", dpi=120)
plt.close()
"/home/claude/har/nbassets/confusion_matrices.png"
\end{verbatim}
\captionof{figure}{\label{orgf026a8f}Confusion matrices (majority-vote cluster to activity mapping)}

\url{Traceback (most recent call last): File "<stdin>", line 1, in <module> File "/tmp/babel-cBWeqx/python-kkGzuk", line 15, in <module> cm\_km, map\_km = cluster\_to\_activity\_confusion(labels\_km) NameError: name 'labels\_km' is not defined}

\begin{verbatim}
print("K-Means cluster -> activity mapping:", {int(k): act_order[v] for k, v in map_km.items()})
print("DBSCAN cluster -> activity mapping:", {int(k): act_order[v] for k, v in map_db.items()})
print("Hierarchical cluster -> activity mapping:", {int(k): act_order[v] for k, v in map_hac.items()})
\end{verbatim}

\begin{verbatim}
Traceback (most recent call last):
  File "<stdin>", line 1, in <module>
  File "/tmp/babel-cBWeqx/python-OgvLKZ", line 1, in <module>
    print("K-Means cluster -> activity mapping:", {int(k): act_order[v] for k, v in map_km.items()})
NameError: name 'map_km' is not defined
\end{verbatim}


\begin{verbatim}
fig, axes = plt.subplots(1, 2, figsize=(12, 4.5))
algos = list(metrics.keys())
xpos = np.arange(len(algos))
axes[0].bar(xpos - 0.2, [metrics[a]["Silhouette"] for a in algos], width=0.4, label="Silhouette")
axes[0].bar(xpos + 0.2, [metrics[a]["Davies_Bouldin"] for a in algos], width=0.4, label="Davies-Bouldin")
axes[0].set_xticks(xpos); axes[0].set_xticklabels(algos, rotation=15); axes[0].legend(); axes[0].set_title("Internal Metrics")
axes[1].bar(xpos - 0.2, [metrics[a]["ARI"] for a in algos], width=0.4, label="ARI")
axes[1].bar(xpos + 0.2, [metrics[a]["NMI"] for a in algos], width=0.4, label="NMI")
axes[1].set_xticks(xpos); axes[1].set_xticklabels(algos, rotation=15); axes[1].legend(); axes[1].set_title("External Metrics (vs Ground Truth)")
plt.tight_layout()
plt.savefig("/home/claude/har/nbassets/metrics_comparison.png", dpi=120)
plt.close()
"/home/claude/har/nbassets/metrics_comparison.png"
\end{verbatim}
\captionof{figure}{\label{org6ad6344}Internal and external evaluation metrics across algorithms}

\begin{center}
\includegraphics[width=.9\linewidth]{Traceback (most recent call last): File "<stdin>", line 1, in <module> File "/tmp/babel-cBWeqx/python-ceaPtG", line 11, in <module> plt.savefig("/home/claude/har/nbassets/metrics_comparison.png", dpi=120) File "/usr/lib/python3/dist-packages/matplotlib/pyplot.py", line 958, in savefig res = fig.savefig(*args, **kwargs) File "/usr/lib/python3/dist-packages/matplotlib/figure.py", line 3019, in savefig self.canvas.print_figure(fname, **kwargs) File "/usr/lib/python3/dist-packages/matplotlib/backend_bases.py", line 2319, in print_figure result = print_method( File "/usr/lib/python3/dist-packages/matplotlib/backend_bases.py", line 1648, in wrapper return func(*args, **kwargs) File "/usr/lib/python3/dist-packages/matplotlib/_api/deprecation.py", line 412, in wrapper return func(*inner_args, **inner_kwargs) File "/usr/lib/python3/dist-packages/matplotlib/backends/backend_agg.py", line 541, in print_png mpl.image.imsave( File "/usr/lib/python3/dist-packages/matplotlib/image.py", line 1675, in imsave image.save(fname, **pil_kwargs) File "/usr/lib/python3/dist-packages/PIL/Image.py", line 2209, in save fp = builtins.open(filename, "w+b") FileNotFoundError: [Errno 2] No such file or directory: '/home/claude/har/nbassets/metrics_comparison.png'}
\end{center}

\section{Discussion}
\label{sec:orgcd6eef8}

\textbf{Which algorithm produced the most meaningful clusters? Why?} By external
agreement with the true activity labels, Hierarchical (Ward) clustering
performed best (ARI = 0.494, NMI = 0.622), narrowly ahead of K-Means
(ARI = 0.420, NMI = 0.560). Both clearly separate static activities
(LAYING/STANDING/SITTING) from dynamic ones (WALKING variants), but
neither cleanly splits SITTING from STANDING or the three WALKING
sub-types -- these have very similar sensor-feature signatures. DBSCAN,
judged by internal cohesion (Silhouette = 0.383, Davies-Bouldin = 0.525 --
both far better than the other two), produced the most internally
compact clusters, but its lower ARI/NMI shows those dense regions don't
align with activity boundaries as well as the partition-based methods.

\textbf{How sensitive was K-Means to the choice of k?} Very sensitive. The
Elbow/Silhouette curves show the data is internally best explained by
\(k=2\) (Silhouette 0.44, essentially "moving" vs. "not moving"),
degrading sharply for every larger \(k\). Forcing \(k=6\) to match the true
activity count drops Silhouette to 0.14 -- the six activities are not
equally well-separated in feature space.

\textbf{Did DBSCAN detect noise or small clusters effectively?} Yes. The eps
sweep showed noise fall smoothly from 46.5\% to under 3\% as eps grew; at
the chosen eps=1.0 it isolated 521 points (5.1\%) as noise -- plausibly
transition frames or atypical movement signatures. It also found small,
tight clusters alongside two large dominant clusters, showing its
strength at picking out small dense sub-populations that K-Means/Ward
tend to absorb into larger neighbours.

\textbf{How does linkage choice (single/complete/ward) affect hierarchical
clustering?} Only Ward's linkage was run here, but its variance-
minimizing objective is similar in spirit to K-Means' WCSS, built up
hierarchically rather than committing to \(k\) centroids up front -- this
is likely why it achieved the best external agreement of the three
algorithms. Single linkage is expected to perform worse here since it is
prone to "chaining" through near-duplicate transitional samples between
activities; complete linkage tends to produce more uneven,
distance-sensitive splits than Ward's variance-based criterion.

\textbf{Which internal metric best matched your visual intuition of cluster
quality?} Silhouette Score matched the 2D visualizations best -- DBSCAN's
much higher Silhouette (0.383) visibly corresponds to its two large,
well-separated blobs, while K-Means/Hierarchical's much lower Silhouette
(\textasciitilde{}0.14) corresponds to the overlapping SITTING/STANDING and
WALKING-variant regions seen throughout. Davies-Bouldin agreed in
direction but was less intuitive to relate to the plots directly;
Calinski-Harabasz favored DBSCAN even more strongly, mostly because it
rewards DBSCAN's tight 2D clusters more than it penalizes the excluded
noise points.

\section{Conclusion}
\label{sec:orge3ef962}

All three algorithms recover the broad "dynamic vs. static" activity
split reliably, but none fully separates all six activities from
unsupervised sensor features alone -- SITTING/STANDING and the three
WALKING variants remain the persistent confusions across K-Means,
DBSCAN, and Hierarchical clustering alike. Hierarchical (Ward) clustering
gave the best agreement with ground truth (ARI 0.494, NMI 0.622), DBSCAN
gave the most internally cohesive clusters while also being the only
algorithm to explicitly flag outliers as noise, and K-Means was most
sensitive to the choice of \(k\), with its own internal metrics actually
favoring a much coarser \(k=2\) split over the six true activity classes.
This highlights a broader lesson: internal cluster-validity metrics
(Silhouette, WCSS) do not necessarily agree with external, label-based
validity -- the two must be interpreted together, not interchangeably.

\section{References}
\label{sec:orgbb791af}

[1] D. Anguita, A. Ghio, L. Oneto, X. Parra, and J. L. Reyes-Ortiz, "A
Public Domain Dataset for Human Activity Recognition Using Smartphones,"
in Proc. 21st Eur. Symp. Artificial Neural Networks (ESANN), Bruges,
Belgium, 2013.

[2] F. Pedregosa et al., "Scikit-learn: Machine Learning in Python,"
Journal of Machine Learning Research, vol. 12, pp. 2825-2830, 2011.

[3] J. MacQueen, "Some methods for classification and analysis of
multivariate observations," in Proc. 5th Berkeley Symp. Mathematical
Statistics and Probability, vol. 1, pp. 281-297, 1967.

[4] M. Ester, H.-P. Kriegel, J. Sander, and X. Xu, "A density-based
algorithm for discovering clusters in large spatial databases with
noise," in Proc. 2nd Int. Conf. Knowledge Discovery and Data Mining
(KDD), pp. 226-231, 1996.

[5] J. H. Ward Jr., "Hierarchical grouping to optimize an objective
function," Journal of the American Statistical Association, vol. 58,
no. 301, pp. 236-244, 1963.
\end{document}
